\chapter*{参赛作品简介}
\addcontentsline{toc}{chapter}{参赛作品简介}

{\renewcommand{\arraystretch}{1.6}%
\begin{tabular}{|>{\centering\arraybackslash}m{0.16\textwidth}|m{0.72\textwidth}|}
\hline
\textbf{参赛作品名称} & BioMed QAgent：面向生物医学研究的多源数据智能采集与分析系统 \\ \hline
\textbf{参赛选题} & 赛道2-方向1A \\ \hline
\textbf{参赛作品简介}（300字以内） & 系统采用“开放式推理与确定性数据处理分权”设计：Qwen 等 LLM 驱动的 Pi Agent 负责理解研究目标、发现候选来源、检查执行路线并提交声明式数据需求；TypeScript Dataset Core 负责解析、规范化、多源整合、质量门禁、人在回路和正式发布。系统不将 Agent 工作区中的临时 CSV 视为任务完成，而以包含数据表、表结构、溯源、审计、校验与内容摘要的正式发布物（追加式发布、读取时摘要重验）作为任务完成的正式证据。支持多种静态数据产品，由数据家族协议表达注册表未覆盖的多表拓扑。来源文件经来源资产、SHA-256、版本证据与定位信息进入证据链；字段映射、单位与低置信度抽取不确定时，可触发可持久化的人在回路请求；任务中断后可在静态执行与检查点范围内恢复，无需模型重解释已确认的数据变换。 \\ \hline
\textbf{参赛作品使用Qwen大模型及AI技术说明} & 本作品未将 Qwen 模型写死在配置中，而是提供 API key 配置方法，允许用户自选包含 Qwen API key 在内的 LLM 作为主 Agent。实验部分（第五章）中，Qwen 系列承担 gold8 正式发布（qwen3.8-27B）、图表理解链路（DashScope Qwen-VL）与全谱系对抗验证（qwen3.7/3.8 各档位）；另以 DeepSeek-V4-Flash 作为评测主力，验证系统的模型无关性。使用 DashScope API key 时，系统可调用 Qwen-VL 理解页面图像进行数据提取。 \\ \hline
\textbf{宣传视频链接/其他作品相关文件链接} & \url{https://github.com/Lidozs55/BioMed-QAgent} \\ \hline
\end{tabular}}
